\subsection{A global arithmetic splitting cannot align two odd Frobenius groups}

\begin{result}{E.186}{Proved}{The sharp one-prime limit on an untwisted arithmetic splitting}
Let
\[
 G=G_{\mathbb Q}(2),\qquad
 \Gamma=G_{\mathbb Q,\{2,\infty\}}(2),\qquad
                  \pi:G\twoheadrightarrow\Gamma.
\]
Fix an actual complex conjugation \(c\in G\), and write
\(\bar c=\pi(c)\). Suppose \(s:\Gamma\to G\) is a
continuous splitting with \(s(\bar c)=c\).
For an odd prime \(q\), let \(f_q\) be an arithmetic
Frobenius in the unramified local decomposition group
of \(\Gamma\). Call \(q\) compatible with \(s\) if
\(s(f_q)\), up to conjugacy in \(G\), is an actual
local Frobenius lift at \(q\).

At most one odd prime is compatible with \(s\), and any
such prime satisfies \(q\equiv3\) or \(5\pmod8\).
Conversely, for any one prime in those two residue classes,
there exists such a splitting compatible with it. Thus
the bound is sharp.

In particular no single real-preserving splitting of
\(\pi\) simultaneously identifies two odd local Frobenius
subgroups with their actual arithmetic lifts. This excludes
that way of removing the path twists in an arithmetic
free-inertia construction. It does not exclude an equivariant
map that retains those twists, and is not a finite-support
theorem for ordinary section pairs.
\end{result}

\begin{proof}
Every quadratic character of \(G\) is a Kummer character.
The characters of \(\Gamma\) have basis
\(\kappa_2(-1),\kappa_2(2)\), as in E.185.6.
For each positive odd prime \(r\), the character
\(\kappa_2(r)\circ s\) vanishes at \(\bar c\): the
actual complex conjugation fixes the real quadratic field
\(\mathbb Q(\sqrt r)\). Consequently there is a uniquely
determined \(e_r\in\mathbb F_2\) such that
\[
                   \kappa_2(r)\circ s=e_r\kappa_2(2).
 \tag{E.186.1}
\]
This uses the real condition on the splitting, not merely
the abstract free-product isomorphism for \(\Gamma\).

Write \(\epsilon(u/q)\in\mathbb F_2\) for the quadratic
residue character, so that
\((u/q)=(-1)^{\epsilon(u/q)}\). If \(q\) is compatible,
then, for every positive odd prime \(r\ne q\), the
unramified character \(\kappa_2(r)\) has its usual
Frobenius value. Conjugation and changing a Frobenius lift
by inertia at \(q\) do not affect that value. Equation
E.186.1 therefore gives
\[
                 \epsilon(r/q)=e_r\epsilon(2/q)
                            \quad(r\ne2,q).
 \tag{E.186.2}
\]
If \(\epsilon(2/q)=0\), every positive odd prime other
than \(q\) would be a square modulo \(q\). Two is also
a square. Choose a positive odd integer coprime to \(q\)
that is a nonsquare modulo \(q\), and factor it into
positive primes. Every prime factor would have square
residue, a contradiction. Thus \((2/q)=-1\), which is
equivalent to \(q\equiv3\) or \(5\pmod8\).

Suppose two distinct compatible primes \(q_1,q_2\) existed.
For every positive odd prime \(r\notin\{q_1,q_2\}\),
equation E.186.2 would imply
\[
                         (r/q_1)=(r/q_2).
 \tag{E.186.3}
\]
Choose by the Chinese remainder theorem a positive odd
integer \(u\), coprime to \(q_1q_2\), with
\((u/q_1)=-1\) and \((u/q_2)=1\).
Its prime factors avoid \(2,q_1,q_2\), and multiplicativity
of E.186.3 would give \((u/q_1)=(u/q_2)\). This is the
required contradiction. The argument uses no density
assertion or infinite Chebotarev intersection.

For sharpness, fix \(q\equiv3\) or \(5\pmod8\) and
choose an actual local Frobenius lift \(\sigma_q\in G\).
The elements \(\bar c\) and \(f_q=\pi(\sigma_q)\)
form a Frattini basis of \(\Gamma\): their evaluations
under \(\kappa_2(-1),\kappa_2(2)\) form a triangular
matrix with diagonal entries one. By E.185,
\(\Gamma\simeq C_2\amalg_{\mathrm{pro}\text{-}2}\mathbb Z_2\).
The homomorphism from this free product sending its
generators to \(\bar c,f_q\) is surjective, hence an
isomorphism by the Hopf property of finitely generated
profinite groups. Sending the same two free-product
generators to \(c,\sigma_q\) now defines a homomorphism
to \(G\); composing with the inverse isomorphism gives
the required splitting. It retains the actual complex
conjugation and the chosen actual local Frobenius lift.

The obstruction can also be seen directly in the path
terms. Put \(N=\ker\pi\), choose compatible representatives
\(\pi(\sigma_q)=f_q\), and define
\[
                        t_q=\sigma_q s(f_q)^{-1}\in N.
\]
For \(r\ne2,q\), its quadratic character is
\[
              \kappa_2(r)(t_q)
                    =\epsilon(r/q)+e_r\epsilon(2/q).
 \tag{E.186.4}
\]
If the splitting is compatible with one prime \(q_0\),
then for every other odd \(q\) at least one coordinate
in E.186.4 with \(r\notin\{2,q,q_0\}\) is nonzero.
For \((2/q)=1\), choose a nonsquare modulo \(q\) that
is coprime to \(2q q_0\); for \((2/q)=-1\), use the
two-prime Chinese-remainder argument just given.
These coordinates are unchanged by local inertia
adjustments at \(q\), or by conjugation, because they
are global quadratic characters unramified at \(q\).
Thus the failure to align the second Frobenius is already
visible in a quadratic quotient of the free-inertia kernel;
it cannot be removed by changing those frames.
\end{proof}
